\documentclass[english,twoside, openright, hidelinks, 11 pt, class=report,crop=false]{standalone}
\input{../../preamb}
\input{../bm}

\begin{document}
	\newpage
	\section{The Equality Sign, Sets, and Number Lines}
	\index{number}
	\subsection*{The Equality Sign}
	As the name suggests, the \outl{equality sign} \index{equality sign}
	%id{symeq}
	\sym{$ = $}
	%\id
	indicates that something is equal. To what extent and when something can be considered equal is a philosophical discussion, so initially, we are left with this: \textsl{The kind of equality
		%id{symeq}
		\sym{$ = $}
		%\id
		referred to must be understood from the context in which the sign is used}. This way, we can study some basic properties of our numbers and later return to more precise meanings of the sign. \regv
	\spr{
		Common ways to say
		%id{symeq}
		\sym{$ = $}
		%\id
		are:
		\begin{itemize}
			\item ''is equal to'' \\
			\item ''is the same as''\\
			\item ''corresponds to''
		\end{itemize}
	}
	\subsection*{Sets and Number Lines}
	In this book, we will use two ways to represent numbers: numbers as a \textsl{set} and numbers as a \textsl{position on a line}. All representations of numbers are based on the understanding of the numbers 0 and 1.
	\newpage
	%span{numbers_as_amounts}
	\subsubsection*{Numbers as Quantities}
	When we talk about a set, the number 0 means
	%id{0interpretations}
	\footnote{In \refkap{Operations} we will see that there are other interpretations of 0.}
	%\id
	''nothing''. A figure where nothing is present will thus be the same as 0:
	%id{equalszero}
	\[ =0 \]
	%\id
	1 will be represented as a single square:
	\fig{rut1}
	Other numbers will then be defined based on how many unit squares ('ones') we have: \vs
	\fig{rut2}
	\fig{rut3}
	%\span
	%span{numbers_on_numberline}
	\subsubsection*{Numbers as Positions on a Line}
	When we place numbers on a line, 0 will be our starting point:
	\fig{lin0a}
	Then we place 1 a certain distance to the right of 0:
	\fig{lin0b}
	Other numbers will then be defined based on how many unit lengths ('ones') we are away from 0:
	\fig{lin1}
	%\span
	\subsection*{Positive Integers}
	We will soon see that numbers do not necessarily have to be \textsl{whole} numbers of ones, but numbers that \textsl{are} have a special name:\regv
	
	\reg[Positive Integers]{
		Numbers that are a whole number of ones are called
		%id{posintoutl}
		\outl{positive\footnotemark integers}\index{positive integers}.
		%\id
		The positive integers are
		%id{posint}
		\[ 1, 2, 3, 4, 5 \text{ and so on.} \]
		%\id
		Positive integers are also called \outl{natural numbers}\index{numbers!natural}.
	}
	\info{What about 0?}{
		Some authors also include 0 in the concept of natural numbers. In some contexts, this makes sense, while in others it does not.
	}
	%id{posinterpre}
	\footnotetext{What the word \textit{positive} means will be discussed in \refkap{Negative Numbers}.}
	%\id
	\subsubsection*{Inequality Signs}
	We also have symbols that indicate that values are \textsl{not} equal. Sometimes it is enough to say that they are not equal, while other times there is also a need to specify which value is larger (and which is smaller).
	\spr{ \vs
		%id{unequals}
		\begin{center}
			\renewcommand{\arraystretch}{1.2}
			\begin{tabular}{@{}cp{0.4cm}l}
				$\neq$&& ''is not equal to''\\
				$ < $ && ''is less than'' \\
				$ > $ && ''is greater than'' \\
				$ \leq $ && ''is less than or equal to'' \\
				$ \geq $ && ''is greater than or equal to''
			\end{tabular}
			\renewcommand{\arraystretch}{1}
		\end{center}
		%\id
	}
	\newpage
	\eks{ \vsb
		%id{unequalexample}
		\alg{
			0 &\neq 1 \vn
			7&<9 \vn
			5&>2 \vn
			\text{positive even numbers}&\geq 2\\
			3&\leq \text{positive odd numbers}
		}
		%\id
	}
	\info{Reading Direction}{
		Strictly speaking, we only have three symbols for inequality, namely
		%id{neq}
		\sym{$\neq$}
		%\id
		\!\!,
		%id{<}
		\sym{$<$}
		%\id
		and
		%id{leq}
		\sym{$\leq$}
		%\id
		\!\!. The last two symbols require a reading direction from left to right. The inequality
		%id{3<5}
		${3<5}$
		%\id
		is read as ''3 is less than 5''. But we can also say that
		%id{<}
		\sym{$<$}
		%\id
		is a directional symbol, where the pointed end should point toward the smaller value. This means that we can read
		%id{3<5}
		${3<5}$
		%\id
		from the right as ''5 is greater than 3''. The fact that we also use the symbols
		%id{>}
		\sym{$>$}
		%\id
		and
		%id{geq}
		\sym{$\geq$}
		%\id
		is a result of sometimes wanting to say ''greater than'' when reading from the left.
	}
	
	\newpage
	%span{integers_larger_than_9}
	\section{Numbers, Digits, and Value \label{numdigval}}
	Our numbers are built from the \outl{digits}\index{digits}
	0, 1, 2, 3, 4, 5, 6, 7, 8, and 9, and their \textsl{position}. The digits and their positions define
	%id{10pos}
	\footnote{Eventually, we will also see that \textit{sign} helps define the value of the number (see \refkap{Negative Numbers}).}
	%\id
	the \outl{value} \index{value} of the number.
	
	\subsection*{Integers Greater Than 9}
	Let's take the number 'fourteen' as an example and write it using our digits.
	\fig{tel_eng}
	A group of ten 'ones' we call a 'ten'. From 'fourteen', we can make 1 'ten', and in addition, we have 4 'ones'. So we write 'fourteen' like this:
	%id{14}
	\[ \text{fourteen}=14 \]
	%\id
	\amounts{\fig{tel2_eng}}
	\numrline{\fig{tel2t}}
	\vsk
	%/span
	\newpage
	\subsection*{Decimal Numbers}
	In many cases, we do not have a whole number of ones, and then we will need to divide 1 into smaller parts. Let's start by drawing a
	%id{defof1}
	\footnote{Here we define the unit length as larger than before.}
	:\!\! 
	%\id
	\amounts{\fig{maal} \vs}
	\numrline{\fig{des1}}
	Then we divide our unit into 10 smaller parts:
	\amounts{\fig{maal1} \vs} \vs
	\numrline{\fig{des1a}}
	Since we have divided 1 into 10 parts, we call such a part a 'tenth':
	\amounts{\fig{maal1a_bm} \vs}
	\numrline{\fig{des1b_bm} \vs}
	
	Tenths are written using the \outl{decimal point} \sym{.}:
	\amounts{\fig{maal1b_eng} \vs}
	\numrline{\fig{des1c_eng} \vs}
	\eks[]{\vs
		\fig{maal2a_eng}
		\fig{des3_eng}
	}
	\regv
	\spr{
		In many countries, a comma \sym{,} is used as the decimal separator:
		%id{dotcomma}
		\alg{
			3,5&\quad(Norwegian, German\;and\;` others) \\
			3.5&\quad (English)
		}
		%\id
		When we combine mathematical symbols
		%id{digits}
		\footnote{Digits are considered mathematical symbols.}
		%\id
		\!\!, we get what we call a mathematical \outl{expression}.
	}
	
	\newpage
	%span{the_decimal_system}
	\subsection*{The Decimal System}
	We have now seen how we can express the value of numbers by placing digits according to the number of tens, ones, and tenths, and it doesn't stop there: \regv
	
	\reg[The Decimal System \label{decimalsys}]{
		The value of a number is given by the digits 0, 1, 2, 3, 4, 5, 6, 7, 8, and 9, and their position. With the digit representing ones as the starting point:
		\begin{itemize}
			\item digits to the left (in order) represent the number of tens, hundreds, thousands, and so on.
			\item digits to the right (in order) represent the number of tenths, hundredths, thousandths, and so on.
		\end{itemize}
	}
	\eks[1]{\vs
		\fig{maal3_eng}
	}
	\eks[2]{ \vs \vs
		\fig{titalsys_eng}
	}
	%span
	\newpage
	\reg[Even and Odd Numbers \label{evenodd}]{
		Integers that have 0, 2, 4, 6, or 8 in the ones place are called \outl{even numbers}\index{even numbers}.\vsk
		
		Integers that have 1, 3, 5, 7, or 9 in the ones place are called \outl{odd numbers}\index{odd numbers}.
	}
	\eks{
		The first ten (positive) even numbers are
		%id{even2to18}
		\[ \text{0, 2, 4, 6, 8, 10, 12, 14, 16, and 18} \]
		%\id
		The first ten (positive) odd numbers are
		%id{odd1to19}
		\[ \text{1, 3, 5, 7, 9, 11, 13, 15, 17, and 19} \]
		%\id
	}
	
	\newpage
	\section{Coordinate System \label{Coord}}
	In many cases, it is useful to use two number lines simultaneously. This is called a \outl{coordinate system}\index{coordinate system}
	%id{coord}
	\!\!\footnote{Strictly speaking, there are many types of coordinate systems, but in this book, we use the term for only one type, namely the \outl{Cartesian coordinate system}. It is named after the French philosopher and mathematician René Descartes.}
	%\id
	\!\!. We then place one number line (an axis) that runs \textsl{horizontally} and one that runs \textsl{vertically}. A position in a coordinate system is called a \outl{point}\index{point}.
	\regdef[Coordinate System]{A point is written as two numbers inside a parenthesis. The two numbers are called the \outl{first coordinate} and the \outl{second coordinate} of the point. \os
		
		The first coordinate shows the point's position along the horizontal axis.\os
		
		The second coordinate shows the point's position along the vertical axis.
	}
	\eks[1]{
		In the figure, we see the points
		%id{p23}
		$ (2, 3) $
		%\id
		\!\!,
		%id{p51}
		$ (5, 1) $
		%\id
		and
		%id{p00}
		$ (0, 0) $
		%\id
		\!\!.
		\fig{kord_eng}
	}
	\spr{
		The point where the axes meet, that is,
		%id{p00comma}
		$ (0, 0) $
		%\id
		\!\!, is called the \outl{origin}\index{origin}.
	}
	\newpage
	\info{Note}{
		The size of a unit length can be different on the two axes.
	}
	\eks[2]{
		\fig{kord2}
	}
\end{document}